\documentclass[aspectratio=169]{beamer}

\usepackage{fontawesome5}
\usepackage{styles/couleurs}
\usepackage{styles/maths}

\usepackage{amsmath}
\usepackage{amssymb}
\usepackage{tikz}

%==================================================
% Style général
%==================================================

\setbeamertemplate{navigation symbols}{}
\setbeamertemplate{footline}[frame number]

\setbeamercolor{frametitle}{
  fg=SavoirC,
  bg=white
}

\setbeamerfont{frametitle}{
  series=\bfseries,
  size=\Large
}

\setbeamerfont{normal text}{
  size=\large
}

\newcommand{\consigne}{%
\begin{beamercolorbox}[
  rounded=true,
  sep=0.8em,
  wd=\textwidth
]{block body}
\centering
\textbf{1. Trouve l'opération qui traduit la situation.}

\medskip

\textbf{2. Calcule le résultat.}
\end{beamercolorbox}
}

%==================================================
% DOCUMENT
%==================================================

\begin{document}

%--------------------------------------------------
\begin{frame}
\frametitle{Réactiver additions et soustractions}

\centering

{\Huge\bfseries Nombres relatifs}

\vspace{0.6cm}

{\Large
Pour chaque situation :

\medskip

\textbf{trouver l'opération, puis calculer.}
}

\end{frame}

%--------------------------------------------------
\begin{frame}
\frametitle{Situation 1 -- Kelvin et degrés Celsius}

Le \textbf{kelvin}, noté \(\mathrm{K}\), est une unité de température.

\medskip

Le zéro absolu correspond à \(0\,\mathrm{K}\).

\medskip

Pour passer d'une température en kelvins à une température
en degrés Celsius, on \textbf{soustrait 273}.

\vspace{0.6cm}

\begin{block}{Question}
À combien de degrés Celsius correspond \(0\,\mathrm{K}\) ?
\end{block}

\vfill

\consigne

\end{frame}

%--------------------------------------------------
\begin{frame}
\frametitle{Situation 2 -- Du Celsius au kelvin}

Le record absolu de froid en France a été enregistré
à une température de

\[
-41\,^\circ\mathrm{C}.
\]

\vspace{0.5cm}

Pour passer des degrés Celsius aux kelvins,
on effectue l'opération inverse.

\vspace{0.4cm}

\begin{block}{Question}
À quelle température en kelvins correspond
\(-41\,^\circ\mathrm{C}\) ?
\end{block}

\vfill

\consigne

\end{frame}

%--------------------------------------------------
\begin{frame}
\frametitle{Situation 3 -- Températures sur la Lune}

Sur la Lune, il peut faire :

\[
130\,^\circ\mathrm{C}
\]

dans les zones exposées au Soleil,

\medskip

et

\[
-150\,^\circ\mathrm{C}
\]

dans les zones à l'ombre.

\vspace{0.4cm}

\begin{block}{Question}
Quel est l'écart entre ces deux températures ?
\end{block}

\vfill

\consigne

\end{frame}

%--------------------------------------------------
\begin{frame}
\frametitle{Situation 4 -- Une matinée froide}

À \(7\) h, la température est de

\[
-6\,^\circ\mathrm{C}.
\]

À midi, elle est de

\[
4\,^\circ\mathrm{C}.
\]

\vspace{0.4cm}

\begin{block}{Question}
De combien de degrés la température a-t-elle augmenté ?
\end{block}

\vfill

\consigne

\end{frame}

%--------------------------------------------------
\begin{frame}
\frametitle{Situation 5 -- Le sous-marin}

Un sous-marin se trouve à l'altitude

\[
-35\text{ m}.
\]

Il \textbf{remonte de \(18\) m}.

\vspace{0.5cm}

\begin{block}{Question}
À quelle altitude se trouve-t-il maintenant ?
\end{block}

\vfill

\consigne

\end{frame}

%--------------------------------------------------
\begin{frame}
\frametitle{Situation 6 -- Le sous-marin continue}

Le sous-marin se trouve maintenant à l'altitude obtenue
dans la situation précédente.

\medskip

Il \textbf{descend ensuite de \(27\) m}.

\vspace{0.5cm}

\begin{block}{Question}
À quelle altitude se trouve-t-il maintenant ?
\end{block}

\vfill

\consigne

\end{frame}

\end{document}